\section{Постановка задачи}
\CWProblem{Разработать библиотеку B-дерева и программу-словарь на её основе.
Ключом является регистронезависимое слово из английских букв длиной не более
256 символов. Значение --- беззнаковое 64-битное целое от $0$ до $2^{64}-1$.
Разные слова могут иметь одинаковые значения.}

\textbf{Вариант:} \VariantNumber, B-дерево.

\begin{tabularx}{\linewidth}{@{}p{35mm}Y@{}}
\hline
Команда & Требуемый результат \\
\hline
\texttt{+ word value} & Добавление: \texttt{OK} или \texttt{Exist}; значение существующего ключа не меняется. \\
\texttt{- word} & Удаление: \texttt{OK} или \texttt{NoSuchWord}. \\
\texttt{word} & Поиск: \texttt{OK: value} или \texttt{NoSuchWord}. \\
\texttt{! Save path} & Сохранение в компактном двоичном виде. \\
\texttt{! Load path} & Загрузка с полной заменой рабочего словаря только после успеха. \\
\hline
\end{tabularx}

Ввод обрабатывается до EOF. При системной ошибке выводится английская
диагностика с префиксом \texttt{ERROR:}. Пустой файл и отсутствие файла
при загрузке считаются допустимыми: результат --- \texttt{OK} и пустой словарь.
При ошибке чтения или неверном формате непустого файла рабочий словарь
сохраняется. Ограничения: 15 секунд и 512 МБ.

\subsection{Исследование качества программы}
По заданию G исследовать время и потребление оперативной памяти,
исправить найденные ошибки или явные недочёты. Отчёт должен содержать
дневник действий, использованные средства и результаты этапов,
выводы о недочётах, сравнение исправленной и предыдущей версий, общие выводы.

В условии названы gprof и dmalloc; разрешены аналогичные или более развитые
средства. Реальные инструменты и версии указать в методике после выполнения.
Основа оформления --- шаблон \cite{reportpattern}.
